\section{Questions about a fixed graph and its colorings}
\label{sec:5.1}
In this chapter we impose the road-coloring hypotheses stated at the
opening. The normalized matrix $A=(R+B)/2$ and its stationary row $\pi$
are fixed while the outgoing red/blue choices vary. Three questions must
be distinguished: whether a CC coloring exists; what kernel rank a given
coloring has; and whether a linear-algebraic test certifies synchronization.

The target-pair graph $H_A$ is defined in
Definition~\ref{def:target-pair-graph}. It is unchanged by recoloring.
We first relate its components to the ranks of $A$ and $\Delta$, then
use bipartiteness to characterize CC existence. We next study consequences
for kernel rank and the vector spaces generated by the semigroup.
Sufficient ranked tests are developed in the following rank-test chapter.
The full sixteen-row four-state run is collected in
Section~\ref{app:four-state-recolorings}, while the CC construction itself
is illustrated locally in Figure~\ref{fig:cc-four-state-target-pair}.

\begin{definition}\label{defn:5.1}
For a coloring semigroup $S=\langle R,B\rangle$, the kernel $K$ is its
minimal nonempty two-sided ideal, equivalently its minimum-rank
transformations. Their common rank is the kernel rank. A synchronizing
instruction is a word whose transformation has rank one.
\end{definition}
